% =====================================================================
\section{Shading and Lighting}\label{sec:shading}
\subsection{Flat, Gouraud and Phong: where lighting is evaluated}
Flat, Gouraud and Phong differ in \emph{where} the lighting model is evaluated, as the proposal requires.
\emph{Flat} (the underground bunker, apron buildings): a geometry shader computes one normal per triangle,
$\vect N_f=\mathrm{normalize}((\vect p_1-\vect p_0)\times(\vect p_2-\vect p_0))$, and lights the triangle once at its
centroid. \emph{Gouraud} (terrain, vegetation, props, base props): lighting in the vertex shader, interpolated.
\emph{Phong} (UAV, rover, tank, helicopter, AA gun, radar): the normal is interpolated and lighting is evaluated per
fragment. Each light's diffuse and specular terms are kept separate across the stages, so shadows and ambient occlusion
(visibility, not lighting) can still be applied per fragment. Test T19 verifies that every leaf has the shading the
proposal prescribes; the keys \texttt{F}, \texttt{G}, \texttt{P}, \texttt{O} force all parts to Flat, Gouraud, Phong or
back to per-object (captures \code{bunker\_flat\_visible}, \code{bunker\_all\_gouraud}, \code{bunker\_all\_phong}).

\begin{figure}[h]
  \centering
  \begin{subfigure}{0.325\linewidth}\pictwo{bunker_flat_visible}{\linewidth}\caption{Flat}\end{subfigure}
  \begin{subfigure}{0.325\linewidth}\pictwo{bunker_all_gouraud}{\linewidth}\caption{Gouraud}\end{subfigure}
  \begin{subfigure}{0.325\linewidth}\pictwo{bunker_all_phong}{\linewidth}\caption{Phong}\end{subfigure}
  \caption{The same object under the three evaluation points (v2 captures of the surface bunker; the v4 base
  uses the same three shaders, with new per-object assignments listed in Table~\ref{tab:parts}).}
  \label{fig:shading}
\end{figure}

\subsection{Cook--Torrance (GGX) lighting}
The visible view uses the Cook--Torrance microfacet BRDF with roughness $\alpha=\rho^2$ and metalness $m$ per material:
\begin{equation}
  f_s=\frac{D\,G\,F}{4(\vect N\!\cdot\!\vect V)(\vect N\!\cdot\!\vect L)},\quad
  D=\frac{\alpha^2}{\pi\big((\vect N\!\cdot\!\vect H)^2(\alpha^2-1)+1\big)^2},\quad
  F=F_0+(1-F_0)(1-\vect V\!\cdot\!\vect H)^5,
\end{equation}
with Smith--Schlick $G$, $F_0=\mathrm{mix}(0.04,\text{albedo},m)$ and diffuse $(1-F)(1-m)$. Key \texttt{B} switches back
to the classic Blinn--Phong $k_d(\vect N\!\cdot\!\vect L)+k_s(\vect N\!\cdot\!\vect H)^n$ of the proposal. Ambient light
is a sky/ground hemisphere. v4 materials are night materials (dark albedos, roughly 0.02--0.30 linear, with an explicit roughness and
metalness chosen per part); T39 checks that every leaf has valid PBR parameters in $[0,1]$.

\subsection{The three proposal lights and the practical lights}
\begin{itemize}
  \item \textbf{L0, moon}: directional, colour $(0.55,0.62,0.80)$, intensity 0.95, shadowed by two cascades.
  \item \textbf{L1, beacon}: point light on the beacon mast above the bunker, colour $(1,0.25,0.10)$,
  attenuation $1/(1+0.35\,d+0.45\,d^2)$, intensity pulsing between 0.6 and 2.0 at 1\,Hz (\code{Scene::applyAgentStates}).
  \item \textbf{L2, UAV searchlight}: spot, $14^\circ$ outer and $9^\circ$ inner cone,
  $f=c^2/(1+0.05\,d+0.03\,d^2)$ with $c$ the smooth-stepped cone factor, intensity 5; it has its own 2048$^2$ (EPIC) shadow map
  and a visible volumetric shaft.
  \item \textbf{Practical lights}: a part whose material is built with \code{asLight(material, intensity, range, cone)}
  is both an emissive object and a light; \nEmissive{} leaves of the export have emissive $>0$. The renderer picks the
  \textbf{16 most relevant lights per view} (\code{kMaxPracticals=16}, \code{gatherPracticals}), ranking by distance divided
  by range. A light's windowed falloff is $f=\mathrm{saturate}(1-(d/R)^4)^2/(1+d^2)$ with an optional spot cone.
  \item \textbf{Dynamic explosion lights.} The combat simulation hands the renderer a list of \code{DynamicLight}
  (blast flashes, burning wrecks, rocket motors, muzzle flashes; position, HDR colour, range). They are appended to the same
  candidate list and rank \emph{ahead} of lamps at the same distance (their score is halved), so an explosion lights
  the base within the 16-light budget instead of being dropped.
\end{itemize}
Keys \texttt{1}--\texttt{4} toggle the moon, beacon, searchlight and the practical lights.

\subsection{Indoor lighting mask}
The underground rooms must not be lit by the moon, the sky or the beacon. \code{WalkSpace} supplies up to eight
axis-aligned \emph{indoor volumes} (the bunker has 5, T69); the renderer uploads the first eight
(\code{uIndoorMin/Max}) and the fragment shader computes a smooth \code{indoorMask(P)}. Inside the mask:
ambient and sky irradiance are multiplied by 0.012, the moon, beacon and searchlight terms by $1-\text{mask}$, fog by
$1-0.85\,\text{mask}$, and every practical light carries an \emph{inside} flag, so a lamp in the tunnel does not
light the surface above it and the other way round (lights on the other side of a wall rank last). The renderer also
disables the lens flare when the camera is indoors, because screens and tubes would ghost everywhere
(\code{FrameState::cameraIndoor}).

\subsection{Procedural surfaces}
There are no image textures. The \code{Surface} enum of \code{Material.h} selects a pattern in the fragment shader;
the v3 set (plain, concrete, ground, metal, paint, foliage, bark, rock, wood, sand/woven, rubber, glass) is kept, and
v4 adds seven (Table~\ref{tab:surfaces}), all implemented in \code{ShaderSources.cpp}.

\begin{table}[h]
\centering\small
\caption{The seven v4 procedural surfaces (\code{Surface} values 12--18).}
\label{tab:surfaces}
\begin{tabular}{@{}rllp{9cm}@{}}
\toprule
Id & Name & Used on & Shader pattern\\\midrule
12 & Corrugated & base containers and huts, radar shelter & ribs every 7.6\,cm along world $Y$ on walls, dirt runs, rust at the bottom\\
13 & Hesco & HESCO bastion walls & welded wire grid (7.6\,cm) over geotextile sand bags, porous when wet\\
14 & Hazard & bunker door portal and sill, jersey barriers, signs & yellow/black $45^\circ$ stripes with wear (base colour ignored)\\
15 & Screen & bunker consoles and tally screens, base control panel & dark phosphor, grid, rotating radar sweep, blips and text bars; use with emissive $>0$\\
16 & Asphalt & the gate road & aggregate speckle, cracks, polished wheel lanes\\
17 & Camo & tank, helicopter, AA gun, base camouflage & three-blob woodland camouflage mixed with the base colour\\
18 & SteelPlate & blast doors, tank and AA-gun armour, radar door & 0.6\,m panels, seam lines, rivet rows, worn bright edges\\\bottomrule
\end{tabular}
\end{table}

\paragraph{Ground.} The ground uses a $1024^2$ data map computed once on the CPU from the layout: gravel and
packed dirt on the road, pads, motor pool, depot and the service ring; patchy meadow on the apron; nothing in the
trench strip; the track channel and creek from v3 are kept. The terrain shader mixes turf, compacted track, wet pebbles
and needle litter and adds a per-pixel detail normal; the per-vertex (Gouraud) irradiance $E_v$ is corrected by
$E=E_v\,\frac{\max(\vect N_d\cdot\vect L,0)+\epsilon}{\max(\vect N\cdot\vect L,0)+\epsilon}$ so the lighting stays per vertex
while the relief shows per pixel.

\subsection{Thermal appearance of materials}
Each material has an intrinsic temperature $h_0$ (ground 0.15, bunker 0.34, rover motor box 0.97; a charred wreck starts at
0.95 and cools). The displayed heat is $h=h_0+0.04\min(E,1.5)$ plus low-frequency ground noise, minus $0.05\,w$ when
it rains and the part is passive ($h_0<0.6$). The light term is small and capped on purpose: a long-wave IR camera
hardly sees visible light. The rover leaves heat on the ground: a footprint grid receives $\Delta T=2.2\,|v_\text{wheel}|\,\Delta t$
under each wheel and cools as $T(t)=T_0e^{-t/\tau}$ with $\tau=45$\,s dry and 15\,s soaked (T47--T49); the hottest track (0.51) stays
below the tracker's 0.80 blob threshold. In v4 a burning wreck also heats the ground around it for the 15\,s it burns
(\code{updateHeatTrail}).
